\documentclass{article}

\usepackage{multirow}
\usepackage[table]{xcolor}
\usepackage{amsmath}

% Custom colors
\definecolor{myblue}{RGB}{10,60,200}
\definecolor{softpink}{HTML}{FFB6C1}
\definecolor{mygreen}{RGB}{30,130,80}

\title{Tables and Colors in LaTeX}
\author{Sahil Kumar}
\date{\today}

\begin{document}

\maketitle

%--------------------------------
\section{Basic Tables}
%--------------------------------

The tabular environment is used to create tables.

\begin{table}[h!]
    \centering
    \begin{tabular}{|c|c|l|}
        \hline
        \textbf{Course Code} & \textbf{Semester} & \textbf{Course Name} \\
        \hline
        22CST201 & III & Data Communications \\
        22CST202 & III & Design and Analysis of Algorithms \\
        22CST203 & III & Digital Circuits and Microprocessors \\
        22CST204 & III & Mathematics for AI \\
        22CST205 & III & Object Oriented Analysis and Design \\
        22CST206 & III & Technical Writing \\
        22HST201 & III & Social Sciences and Professional Ethics \\
        \hline
    \end{tabular}
    \caption{Third Semester Courses at MNIT Jaipur}
    \label{tab:courses}
\end{table}

Table~\ref{tab:courses} shows the courses of Semester III.

%--------------------------------
\section{Advanced Tables}
%--------------------------------

\subsection{Multicolumn}

\begin{table}[h!]
    \centering
    \begin{tabular}{|c|c|c|}
        \hline
        \multicolumn{3}{|c|}{\textbf{MNIT Jaipur - Semester III}} \\
        \hline
        Course Code & Semester & Course Name \\
        \hline
        22CST201 & III & Data Communications \\
        22CST202 & III & DAA \\
        \hline
    \end{tabular}
    \caption{Example of Multicolumn}
\end{table}

\subsection{Multirow}

\begin{table}[h!]
    \centering
    \begin{tabular}{|c|c|c|}
        \hline
        \textbf{Semester} & \textbf{Course Code} & \textbf{Course} \\
        \hline
        \multirow{3}{*}{III}
        & 22CST201 & Data Communications \\
        \cline{2-3}
        & 22CST202 & DAA \\
        \cline{2-3}
        & 22CST203 & Digital Circuits \\
        \hline
    \end{tabular}
    \caption{Example of Multirow}
\end{table}

\subsection{Table Positioning}

The table placement options control where LaTeX tries to place a table.

\begin{itemize}
    \item h: Here
    \item t: Top of the page
    \item b: Bottom of the page
    \item p: Separate page for floats
    \item !: Relax placement restrictions
\end{itemize}

%--------------------------------
\section{Colors in LaTeX}
%--------------------------------

\subsection{Basic Text Colors}

\textcolor{red}{Red Text}

\textcolor{blue}{Blue Text}

\textcolor{green}{Green Text}

\textcolor{violet}{Violet Text}

\textcolor{myblue}{Custom Blue Text}

\subsection{Defining Custom Colors}

Custom colors can be defined using RGB and HTML values.

\textcolor{softpink}{This text uses an HTML-defined color.}

\textcolor{mygreen}{This text uses an RGB-defined color.}

\subsection{Colored Mathematical Expressions}

\[
\textcolor{magenta}{E=mc^2}
\]

\[
\textcolor{blue}{\sum_{i=1}^{n} i}
\]

\[
\textcolor{orange}{\int_0^1 x^2\,dx}
\]

%--------------------------------
\section{Colored Tables}
%--------------------------------

The xcolor package with the table option allows
us to add colors to table rows and cells.

\subsection{Colored Rows}

\begin{table}[h!]
    \centering
    \begin{tabular}{|c|c|l|}
        \hline
        \rowcolor{myblue}
        \textcolor{white}{\textbf{Course Code}} &
        \textcolor{white}{\textbf{Semester}} &
        \textcolor{white}{\textbf{Course Name}} \\
        \hline

        \rowcolor{blue!10}
        22CST201 & III & Data Communications \\
        \hline

        \rowcolor{green!15}
        22CST202 & III & Design and Analysis of Algorithms \\
        \hline

        \rowcolor{yellow!20}
        22CST203 & III & Digital Circuits and Microprocessors \\
        \hline

        22CST204 & III & Mathematics for AI \\
        \hline
    \end{tabular}
    \caption{Table with Colored Rows}
\end{table}

\subsection{Colored Cells}

\begin{table}[h!]
    \centering
    \begin{tabular}{|c|c|c|}
        \hline
        Course & Semester & Status \\
        \hline
        Data Communications &
        III &
        \cellcolor{green!25} Completed \\
        \hline
        DAA &
        III &
        \cellcolor{yellow!30} Ongoing \\
        \hline
        Digital Circuits &
        III &
        \cellcolor{red!20} Pending \\
        \hline
    \end{tabular}
    \caption{Table with Colored Cells}
\end{table}

\end{document}